%%=============================================================================
%% Inleiding
%%=============================================================================

\chapter{\IfLanguageName{dutch}{Inleiding}{Introduction}}%
\label{ch:inleiding}

In moderne zorglabs groeit het aantal verbon den apparaten snel. Denk dan maar aan IoT apparatuur die continu data uitwisselen. De hui dige WiFi-infrastructuren botsen echter op be perkingen zoals interferentie, beperkte dekking en uitdagingen rond schaalbaarheid en beveiliging.
\vspace{0.5cm}

5G biedt een potentiële oplossing door hogere bandbreedte, lagere latentie en betere ondersteuning voor IoT-communicatie. Bovendien maakt 5G een hogere mate van mobiliteit en netwerksegmentatie mogelijk. In het zorg lab komt daarom de vraag of een private 5G-netwerkomgeving een geschikt en veilig alternatief kan vormen voor de bestaande infrastructuren.
\vspace{0.5cm}

Tegelijkertijd brengt 5G nieuwe risico’s met zich mee. Voorbeelden hiervan zijn kwetsbaarheden binnen slicing-architecturen, configuratiefouten in private 5G-netwerken, mogelijke interceptie van gevoelige gezondheidsdata en de noodzaak van streng device- en SIM-beheer. Voor zorginstellingen, waar vertrouwelijkheid, integriteit en beschikbaarheid van data essentieel zijn, is het noodzakelijk deze risico’s grondig te analyseren en gepaste beveiligingsmaatregelen te definiëren.


\section{\IfLanguageName{dutch}{Probleemstelling}{Problem Statement}}%
\label{sec:probleemstelling}

Zorglabs worden dagelijks geconfronteerd met toenemende connectiviteitsvereisten door de groei van zorg  IoT-toestellen, terwijl bestaande WiFi-infrastructuren beperkingen hebben op vlak van schaalbaarheid, betrouwbaarheid, beveiliging en snelheid. Hoewel private 5G een alternatief vormt, is het nog steeds onvoldoende duidelijk hoe deze technologie veilig en haalbaar kan worden geïmplementeerd binnen een zorgomgeving met strikte eisen rond beveiliging.


\section{\IfLanguageName{dutch}{Onderzoeksvraag}{Research question}}%
\label{sec:onderzoeksvraag}

De centrale onderzoeksvraag van dit onderzoek is als volgt:
Hoe kan een private 5G-netwerkomgeving veilig en technisch haalbaar worden opgezet en geconfigureerd voor zorglabtoestellen, en welke beveiligingsmaatregelen zijn hierbij noodzakelijk?


\section{\IfLanguageName{dutch}{Onderzoeksdoelstelling}{Research objective}}%
\label{sec:onderzoeksdoelstelling}

Het doel van dit onderzoek is het uitwerken van een Proof of concept voor private 5G binnen een zorglabomgeving. De resultaten richten zich op IT-beheerders, cybersecurity-specialisten en onderzoekers binnen de zorgsector die veilige en toekomstbestendige 5G-connectiviteit willen realiseren.

\section{\IfLanguageName{dutch}{Opzet van deze bachelorproef}{Structure of this bachelor thesis}}%
\label{sec:opzet-bachelorproef}

% Het is gebruikelijk aan het einde van de inleiding een overzicht te
% geven van de opbouw van de rest van de tekst. Deze sectie bevat al een aanzet
% die je kan aanvullen/aanpassen in functie van je eigen tekst.

De rest van deze bachelorproef is als volgt opgebouwd:

In Hoofdstuk~\ref{ch:stand-van-zaken} wordt een overzicht gegeven van de stand van zaken binnen het onderzoeksdomein, op basis van een literatuurstudie.

In Hoofdstuk~\ref{ch:methodologie} wordt de methodologie toegelicht en worden de gebruikte onderzoekstechnieken besproken om een antwoord te kunnen formuleren op de onderzoeksvragen.

% TODO: Vul hier aan voor je eigen hoofstukken, één of twee zinnen per hoofdstuk

In Hoofdstuk~\ref{ch:conclusie}, tenslotte, wordt de conclusie gegeven en een antwoord geformuleerd op de onderzoeksvragen. Daarbij wordt ook een aanzet gegeven voor toekomstig onderzoek binnen dit domein.